\begin{textsample}{POS Dim 2 – profiled\_gpt – Score 6.00 – profiled\_gpt/t000418\_gpt.txt}  \label{ex:f2_pos_018}
so what I was trying to say is you can actually turn it round a bit, \textbf{rather} \textbf{than} just going in and showing everyone what you know.
If you get other people to talk about what they don't know, or what they think they know and why, you ’re kind of teasing out their misconceptions.
That creates this sort of reversed, pushback style of learning, where you ’re not just pushing your \textbf{own} \textbf{ideas} at them, you ’re drawing theirs out and then examining them together.
It means \textbf{more} people are likely to be included, instead of them just sitting there thinking they can't keep up.
You \textbf{build} a bit of rapport because you ’re treating them as fellow learners \textbf{rather} \textbf{than} an audience.
And it gives you an alternative to just performing your \textbf{own} knowledge all the time, which can be quite isolating and, frankly, a bit dull for you and for them.

% matched lemmas: build, idea, more, own, rather, than
\end{textsample}
